% Necoyoad Architecture Blueprint v8 — CMS Posts/Pages and Widget Composition
\documentclass[11pt,a4paper]{article}
\usepackage[T1]{fontenc}
\usepackage[utf8]{inputenc}
\usepackage{lmodern}
\usepackage{microtype}
\usepackage{graphicx}
\usepackage{xcolor}
\usepackage{geometry}
\usepackage{amsmath}
\usepackage{amssymb}
\usepackage{tcolorbox}
\tcbuselibrary{skins,breakable,listings}
\usepackage{colortbl}
\usepackage{booktabs}
\usepackage{tabularx}
\usepackage{longtable}
\usepackage{array}
\usepackage{enumitem}
\usepackage{adjustbox}
\usepackage{caption}
\usepackage{fancyhdr}
\usepackage{titlesec}
\usepackage{pdflscape}
\usepackage{listings}
\usepackage{tikz}
\usetikzlibrary{arrows.meta,positioning,calc,shapes.geometric,fit,backgrounds,chains,decorations.pathreplacing,shapes.multipart}
\usepackage[colorlinks=true,linkcolor=blue!50!black,citecolor=blue!40!black,urlcolor=blue!60!black,bookmarks=true,bookmarksnumbered=true,unicode=true,pdftitle={Necoyoad - CMS Posts/Pages and Widget Composition (v8)},pdfauthor={Reverse-Engineered Architecture Report},pdfsubject={Software Architecture, CMS, Posts, Pages, Widget Composition, Dynamic Rendering},pdfkeywords={Necoyoad, OpenCart, PHP 8.1, CMS, Posts, Pages, Widget Composition, Dynamic Rendering, Hardcoded Placeholders, Template Override, widgets-common.tpl}]{hyperref}
\geometry{a4paper, top=2.4cm, bottom=2.4cm, left=2.4cm, right=2.4cm}
\setlength{\headheight}{14pt}
\pagestyle{fancy}
\fancyhf{}
\fancyhead[L]{\small\itshape Necoyoad v8 --- CMS Posts/Pages and Widget Composition}
\fancyfoot[C]{\small\thepage}
\renewcommand{\headrulewidth}{0.3pt}
\renewcommand{\footrulewidth}{0pt}
\titleformat{\section}{\Large\bfseries\color{black!85}}{\thesection}{0.6em}{}
\titleformat{\subsection}{\large\bfseries\color{black!80}}{\thesubsection}{0.6em}{}
\titleformat{\subsubsection}{\normalsize\bfseries\color{black!75}}{\thesubsubsection}{0.5em}{}
\widowpenalty=10000
\clubpenalty=10000
\emergencystretch=4em
\hbadness=10000
\hfuzz=2pt
\vfuzz=2pt
\sloppy
\definecolor{necoyoad-blue}{HTML}{1F4E79}
\definecolor{necoyoad-orange}{HTML}{B85C00}
\definecolor{necoyoad-green}{HTML}{2E6B33}
\definecolor{nodefill}{HTML}{F4F7FB}
\definecolor{nodefillalt}{HTML}{FBF4EE}
\definecolor{nodefillgreen}{HTML}{F2F8F2}
\definecolor{nodefillgray}{HTML}{F0F0F2}
\definecolor{phpline}{HTML}{F8F8F0}
\definecolor{phpcomment}{HTML}{8C8C8C}
\definecolor{phpkeyword}{HTML}{1F4E79}
\definecolor{phpstring}{HTML}{B85C00}
\lstdefinelanguage{necophp}{
  morekeywords={abstract,and,array,as,break,case,catch,class,clone,const,continue,declare,default,do,echo,else,elseif,empty,enddeclare,endfor,endforeach,endif,endswitch,endwhile,extends,final,finally,for,foreach,function,global,goto,if,implements,include,include_once,instanceof,insteadof,interface,namespace,new,or,print,private,protected,public,require,require_once,return,static,switch,throw,trait,try,use,var,while,xor,yield,true,false,null},
  sensitive=true,
  morecomment=[l]{//},
  morecomment=[s]{/*}{*/},
  morecomment=[l]{\#},
  morestring=[b]",
  morestring=[b]',
}
\lstset{
  language=necophp,
  basicstyle=\ttfamily\footnotesize,
  keywordstyle=\color{phpkeyword}\bfseries,
  commentstyle=\color{phpcomment}\itshape,
  stringstyle=\color{phpstring},
  backgroundcolor=\color{phpline},
  numbers=left,
  numberstyle=\tiny\color{phpcomment},
  numbersep=8pt,
  frame=single,
  rulecolor=\color{black!25},
  framerule=0.3pt,
  breaklines=true,
  breakatwhitespace=true,
  showstringspaces=false,
  captionpos=b,
  xleftmargin=14pt,
  xrightmargin=4pt,
  aboveskip=8pt,
  belowskip=8pt,
}
\newtcolorbox{findingbox}[1][]{
  enhanced, breakable,
  colback=nodefillalt, colframe=necoyoad-orange!70,
  boxrule=0.4pt, arc=2pt, left=8pt, right=8pt, top=6pt, bottom=6pt,
  fonttitle=\bfseries\small, coltitle=white,
  title=#1
}
\captionsetup{justification=raggedright, singlelinecheck=false, font=small, labelfont=bf}

\begin{document}
\thispagestyle{empty}
\begin{center}
{\Large\bfseries Necoyoad --- CMS Posts/Pages and Widget Composition (v8)}\\[0.4em]
{\small\itshape How posts and pages share a table, how widgets compose with them
dynamically and manually, and how templates can be overridden per entity}\\[1.2em]
\end{center}
\begin{abstract}
\noindent
This is the eighth volume of the Necoyoad architecture blueprint. v8 covers the
CMS subsystem (posts and pages) and its deep integration with the widget
composition system. The key finding is that posts and pages share the same
database table (\texttt{post}) with a \texttt{post\_type} discriminator, and that
every CMS page is a widget composition surface --- the same
\texttt{loadWidgets()}/\texttt{addChild()}/\texttt{render()} pipeline that drives
the home page and product pages also drives blog posts, static pages, and
post categories.

v8 documents two composition modes: \emph{dynamic composition} (the admin
configures widgets via the visual editor, and the \texttt{widgets-rows.tpl}
template emits \texttt{\{\%widget\_name\%\}} placeholders that
\texttt{Controller::fetch()} substitutes at render time) and \emph{manual
hardcoding} (a template author writes \texttt{\{\%widget\_name\%\}} tokens
directly in the \texttt{.tpl} file, bypassing the admin widget manager). It
also covers the per-entity template override mechanism (EAV
\texttt{property('style', 'view')} returns a template path), the
\texttt{widgets-common.tpl} universal layout fragment, the \texttt{only:}
position prefix for embedded pages, and the \texttt{page\_embed.tpl} simplified
layout used by the menu system's \texttt{submenu\_type = 'page\_id'}.

The stance is unchanged: \textbf{descriptive, not prescriptive}.
\end{abstract}
\vspace{0.6em}
\noindent\textbf{Keywords:} Necoyoad, OpenCart, PHP 8.1, CMS, posts, pages,
widget composition, dynamic rendering, hardcoded placeholders, template override,
\texttt{widgets-common.tpl}, \texttt{widgets-rows.tpl}, \texttt{only:} prefix,
\texttt{page\_embed.tpl}, EAV property, polymorphic description.
\newpage
\tableofcontents
\newpage

% ============================================================
\section{Introduction}\label{sec:intro}
% ============================================================

\subsection{Eight Volumes}

v1--v7 covered the whole platform, source verification, the rendering pipeline,
the marketing subsystem, multi-store/multi-language, the admin back-office, and
the menu links system. v8 covers the CMS posts/pages subsystem and its
integration with the widget composition system.

\subsection{Scope}

v8 covers:
\begin{itemize}[leftmargin=1.6em,itemsep=2pt]
  \item The \texttt{post} table --- shared by posts and pages via
        \texttt{post\_type}.
  \item The storefront content controllers (\texttt{content/post},
        \texttt{content/page}, \texttt{content/category}).
  \item The per-entity template override mechanism (EAV
        \texttt{property('style', 'view')}).
  \item The \texttt{widgets-common.tpl} universal layout fragment.
  \item The \texttt{widgets-rows.tpl} placeholder emitter.
  \item Dynamic widget composition (admin-configured, via the visual editor).
  \item Manual widget composition (hardcoded \texttt{\{\%widget\_name\%\}}
        tokens in \texttt{.tpl} files).
  \item The \texttt{only:} position prefix for embedded pages.
  \item The \texttt{page\_embed.tpl} simplified layout.
  \item Landing page templates (combining widget-driven and hardcoded content).
\end{itemize}

% ============================================================
\section{The Shared \texttt{post} Table}\label{sec:post-table}
% ============================================================

\begin{findingbox}[Posts and pages share the same table]
The \texttt{post} table stores both blog posts and static pages. The
\texttt{post\_type} column discriminates them: \texttt{'post'} for blog
posts, \texttt{'page'} for static pages. Both the post model
(\texttt{ModelContentPost}) and the page model
(\texttt{ModelContentPage}) declare \texttt{\$table = "post"} and
\texttt{\$pkey = "post\_id"}. The only difference is
\texttt{\$object\_type}: \texttt{'post'} vs \texttt{'page'}. This
discriminator flows into the polymorphic \texttt{description},
\texttt{property}, and \texttt{object\_to\_store} tables.
\end{findingbox}

\begin{lstlisting}[caption={The \texttt{post} table schema (shared by posts and pages).},label={lst:post-schema}]
CREATE TABLE `post` (
  `post_id`            int(11) NOT NULL,
  `parent_id`          int(11) NOT NULL,          -- tree via parent_id
  `author_id`          int(11) NOT NULL,          -- admin user who wrote it
  `post_type`          varchar(50) NOT NULL,      -- 'post' or 'page'
  `sort_order`         int(11) NOT NULL,
  `image`              varchar(250) NOT NULL,
  `status`             int(1) NOT NULL,
  `date_publish_start` datetime NOT NULL,
  `date_publish_end`   datetime NOT NULL,
  `publish`            int(1) NOT NULL,           -- published flag
  `allow_reviews`      int(1) NOT NULL,
  `template`           varchar(100) NOT NULL,     -- template hint (unused by storefront)
  `date_added`         datetime NOT NULL,
  `date_modified`      datetime NOT NULL
);
\end{lstlisting}

The \texttt{template} column in the \texttt{post} table is declared but
\emph{not used} by the storefront controllers. The actual template override
comes from the EAV \texttt{property} table, not from this column.

% ============================================================
\section{The Storefront CMS Controllers}\label{sec:cms-controllers}
% ============================================================

\subsection{The Common CMS Page Pattern}

All three storefront CMS controllers (post, page, category) follow the
same pattern:

\begin{lstlisting}[caption={The common CMS page pattern (abridged from \texttt{content/post.php}).},label={lst:cms-pattern}]
public function index() {
    $this->language->load('content/post');
    $this->load->model('content/post');

    // 1. Cache key (per-entity + language + currency + store + customer)
    $cacheId = 'html-post.' . $this->request->getQuery('post_id') . ...;

    // 2. Set session object_type/object_id (for per-object widget filtering)
    $this->session->set('object_type', 'post');
    $this->session->set('object_id', $post_id);

    // 3. Load the entity
    $post_info = $this->modelPost->getById($post_id);

    if ($post_info && $post_info['publish']) {
        // 4. Customer group permission check
        $customerGroups = $this->modelPost->getProperty($post_id, 'customer_groups', 'customer_groups');
        if (in_array(0, $customerGroups) || in_array($this->customer->getCustomerGroupId(), $customerGroups)) {

            // 5. Cache check
            $cached = $this->cache->get($cacheId);
            if ($cached && !$this->user->isLogged()) {
                $this->response->setOutput($cached, ...);
            } else {
                // 6. Set landing_page (for widget filtering)
                $this->session->set('landing_page', 'content/post/index');

                // 7. Load widgets for three positions
                $this->loadWidgets('featuredContent');
                $this->loadWidgets('main');
                $this->loadWidgets('featuredFooter');

                // 8. Add layout children
                $this->addChild('common/column_left');
                $this->addChild('common/column_right');
                $this->addChild('common/header');
                $this->addChild('common/footer');

                // 9. Per-entity template override
                $template = $this->modelPost->getProperty($post_id, 'style', 'view');
                $default_template = $this->config->get('default_view_post') ?: 'content/post.tpl';
                $template = empty($template) ? $default_template : $template;
                if (file_exists(DIR_TEMPLATE . $this->config->get('config_template') . '/' . $template)) {
                    $this->template = $this->config->get('config_template') . '/' . $template;
                } else {
                    $this->template = 'choroni/' . $template;
                }

                // 10. Render (with cache write-back)
                if (!$this->user->isLogged()) $this->cacheId = $cacheId;
                $this->response->setOutput($this->render(true), ...);
            }
        }
    } else {
        $this->error404();
    }
}
\end{lstlisting}

\subsection{The Session \texttt{object\_type} Protocol}

Each CMS controller sets the session's \texttt{object\_type} and
\texttt{object\_id}:

\begin{table}[htbp]
\centering
\caption{Session \texttt{object\_type} per CMS controller.}\label{tab:object-types}
\small
\begin{tabularx}{\columnwidth}{llX}
\toprule
\textbf{Controller} & \textbf{object\_type} & \textbf{Effect} \\
\midrule
\texttt{content/post}     & \texttt{'post'}          & Widgets filtered by \texttt{LIKE '\%object\_type=post\%'} and \texttt{LIKE '\%object\_id=<post\_id>\%'}. \\
\texttt{content/page}     & \texttt{'page'}          & Widgets filtered by \texttt{LIKE '\%object\_type=page\%'} and \texttt{LIKE '\%object\_id=<page\_id>\%'}. \\
\texttt{content/category} & \texttt{'post\_category'} & Widgets filtered by \texttt{LIKE '\%object\_type=post\_category\%'} and \texttt{LIKE '\%object\_id=<cat\_id>\%'}. \\
\bottomrule
\end{tabularx}
\end{table}

This is the mechanism v5 documented: \texttt{Controller::loadWidgets()} reads
the session's \texttt{object\_type}/\texttt{object\_id} and runs a second
\texttt{NecoWidget::getRows()} query with \texttt{LIKE} predicates, merging
per-object widgets into the default tree.

\subsection{Customer Group Permission Check}

Every CMS entity has a \texttt{customer\_groups} property in the EAV
\texttt{property} table (under \texttt{group='customer\_groups'},
\texttt{key='customer\_groups'}). This is a serialised array of customer
group IDs that are allowed to view the entity. If the array contains
\texttt{0}, the entity is public (visible to all). Otherwise, the
customer's group must be in the array.

\subsection{The Per-Entity Template Override}

The template resolution follows a 3-level priority:

\begin{enumerate}[leftmargin=1.6em,itemsep=2pt]
  \item \textbf{EAV property} \texttt{property('style', 'view')} --- the
        admin can set a per-entity template path (e.g.\
        \texttt{content/landing\_page\_necotienda.tpl}) via the visual
        editor. This is stored in the \texttt{property} EAV table under
        \texttt{group='style'}, \texttt{key='view'}.
  \item \textbf{Config default} --- \texttt{default\_view\_post} or
        \texttt{default\_view\_page} setting (per-store, in the
        \texttt{setting} table). Falls back to
        \texttt{content/post.tpl} or \texttt{content/page.tpl}.
  \item \textbf{Theme fallback} --- if the template file exists in the
        active theme (\texttt{config\_template}), use it; otherwise fall
        back to \texttt{choroni/}.
\end{enumerate}

\begin{lstlisting}[caption={Per-entity template override (from \texttt{content/post.php}).},label={lst:template-override}]
$template = $this->modelPost->getProperty($post_id, 'style', 'view');
$default_template = ($this->config->get('default_view_post'))
    ? $this->config->get('default_view_post')
    : 'content/post.tpl';
$template = empty($template) ? $default_template : $template;
if (file_exists(DIR_TEMPLATE . $this->config->get('config_template') . '/' . $template)) {
    $this->template = $this->config->get('config_template') . '/' . $template;
} else {
    $this->template = 'choroni/' . $template;
}
\end{lstlisting}

This means each individual post or page can use a completely different
template --- a blog post can use \texttt{content/post.tpl}, a landing page
can use \texttt{content/landing\_page\_necotienda.tpl}, and a product
information page can use a custom \texttt{content/product-info.tpl} --- all
without changing any code.

% ============================================================
\section{The \texttt{widgets-common.tpl} Universal Layout}\label{sec:widgets-common}
% ============================================================

The most important template in the theme is
\texttt{shared/widgets-common.tpl}. It is the universal layout fragment
included by almost every page template:

\begin{lstlisting}[caption={\texttt{shared/widgets-common.tpl} --- the universal layout fragment.},label={lst:widgets-common}]
<?php include("widgets-featured.tpl");?>

<!--mainContentContainer -->
<div id="mainContentContainer" nt-editable>
    <div class="row">

        <?php include("breadcrumbs.tpl");?>
        <div class="clear"></div>

        <!-- left-column -->
        <?php if ($column_left) { ?>
        <?php include("widgets-column-left.tpl");?>
        <?php } ?>
        <!--/left-column -->

        <!--center-column -->
        <?php include("widgets-column-center.tpl");?>
        <!--/center-column -->

        <!-- right-column -->
        <?php if ($column_right) { ?>
        <?php include("widgets-column-right.tpl");?>
        <?php } ?>
        <!--/right-column -->

    </div>
</div>

<?php include("widgets-featured-footer.tpl");?>
\end{lstlisting}

This template includes five fragments:
\begin{enumerate}[leftmargin=1.6em,itemsep=2pt]
  \item \texttt{widgets-featured.tpl} --- sets \texttt{\$position =
        'featuredContent'} and includes \texttt{widgets-rows.tpl}.
  \item \texttt{breadcrumbs.tpl} --- renders the breadcrumb trail.
  \item \texttt{widgets-column-left.tpl} --- echoes \texttt{\$column\_left}
        (the rendered left column child).
  \item \texttt{widgets-column-center.tpl} --- sets
        \texttt{\$position = 'main'} and includes
        \texttt{widgets-rows.tpl}. Also computes the column width based
        on which side columns are present.
  \item \texttt{widgets-column-right.tpl} --- echoes
        \texttt{\$column\_right}.
  \item \texttt{widgets-featured-footer.tpl} --- sets
        \texttt{\$position = 'featuredFooter'} and includes
        \texttt{widgets-rows.tpl}.
\end{enumerate}

\subsection{Which Templates Use \texttt{widgets-common.tpl}}

\begin{table}[htbp]
\centering
\caption{Templates that include \texttt{widgets-common.tpl}.}\label{tab:common-users}
\small
\begin{tabularx}{\columnwidth}{llX}
\toprule
\textbf{Template} & \textbf{Controller} & \textbf{Notes} \\
\midrule
\texttt{content/post.tpl}     & \texttt{content/post}     & Blog post page. \\
\texttt{content/page.tpl}     & \texttt{content/page}     & Static page. \\
\texttt{store/product.tpl}    & \texttt{store/product}    & Product detail page. \\
\texttt{content/category.tpl} & \texttt{content/category} & Blog category page. \\
\texttt{content/categories.tpl} & \texttt{content/category/all} & All blog categories. \\
\texttt{content/posts.tpl}    & \texttt{content/post/all} & All blog posts. \\
\bottomrule
\end{tabularx}
\end{table}

The \texttt{common/home.tpl} (home page) does \emph{not} use
\texttt{widgets-common.tpl} --- it manually includes each fragment
individually. This is an older pattern that predates
\texttt{widgets-common.tpl}. The home page template also omits the
breadcrumbs fragment.

% ============================================================
\section{The \texttt{widgets-rows.tpl} Placeholder Emitter}\label{sec:widgets-rows}
% ============================================================

\texttt{shared/widgets-rows.tpl} is the template that actually emits the
\texttt{\{\%widget\_name\%\}} placeholder tokens. It iterates the
\texttt{\$rows[\$position]} array (populated by
\texttt{Controller::loadWidgets()}) and emits the widget tree as nested
HTML with placeholder tokens:

\begin{lstlisting}[caption={\texttt{shared/widgets-rows.tpl} --- the placeholder emitter.},label={lst:widgets-rows}]
<?php foreach($rows[$position] as $j => $row) { ?>
    <?php if (!$row['key']) continue; ?>
    <?php $row_id = $row['key']; ?>
    <?php $row_settings = unserialize($row['value']); ?>

    <div data-row="<?php echo $row_id; ?>" data-position="<?php echo $position; ?>"
         class="row<?php if (!empty($row_settings['classnames'])) echo ' '.$row_settings['classnames']; ?>"
         id="<?php echo $position; ?>_<?php echo $row_id; ?>" nt-editable>

        <?php foreach($row['columns'] as $k => $column) { ?>
            <?php $column_id = $column['key']; ?>
            <?php $column_settings = unserialize($column['value']); ?>

            <div data-column="<?php echo $column_id; ?>" data-position="<?php echo $position; ?>"
                 class="large-<?php echo $column_settings['grid_large']; ?>
                       medium-<?php echo $column_settings['grid_medium']; ?>
                       small-<?php echo $column_settings['grid_small']; ?>"
                 id="<?php echo $position; ?>_<?php echo $column_id; ?>" nt-editable>
                <ul class="widgets">
                    <?php foreach($column['widgets'] as $l => $widget) { ?>
                        {%<?php echo $widget['name']; ?>%}
                    <?php } ?>
                </ul>
            </div>
        <?php } ?>
    </div>
<?php } ?>
\end{lstlisting}

The key line is:
\texttt{<?php foreach(\$column['widgets'] as \$widget) \{ ?> \{\%<?php echo \$widget['name']; ?>\%\} <?php \} ?>}

This emits \texttt{\{\%widget\_name\%\}} for each widget in each column
in each row. The \texttt{Controller::fetch()} method (v3, \S3.2) then
substitutes each \texttt{\{\%widget\_name\%\}} token with the
corresponding \texttt{\$<name>\_code} output (the widget's rendered HTML).

\subsection{The Row and Column Settings}

Each row and column has a serialised PHP settings blob (in the
\texttt{value} column of the \texttt{property} EAV table). The
\texttt{widgets-rows.tpl} template unserialises these and uses them for:

\begin{itemize}[leftmargin=1.6em,itemsep=2pt]
  \item \textbf{Row settings}: \texttt{sticky} (data-sticky attribute),
        \texttt{layout\_width} (\texttt{'fixed'} adds the
        \texttt{container} class), \texttt{classnames} (additional CSS
        classes).
  \item \textbf{Column settings}: \texttt{grid\_large},
        \texttt{grid\_medium}, \texttt{grid\_small} (Foundation grid
        column widths for responsive layouts), \texttt{sticky},
        \texttt{classnames}.
\end{itemize}

The \texttt{nt-editable} attribute on every row and column div is the
visual theme editor's hook (v3, \S9).

% ============================================================
\section{Dynamic vs Manual Widget Composition}\label{sec:composition-modes}
% ============================================================

Necoyoad supports two ways to compose widgets on a CMS page:
\emph{dynamic} (admin-configured via the visual editor) and \emph{manual}
(hardcoded \texttt{\{\%widget\_name\%\}} tokens in the \texttt{.tpl} file).

\subsection{Dynamic Composition (The Default)}

In dynamic composition, the admin uses the visual theme editor to drag
widgets into rows and columns. The widget tree is stored in the
\texttt{property} EAV table (\texttt{object\_type = 'widget\_rows'},
\texttt{group = <position>}). At render time:

\begin{enumerate}[leftmargin=1.6em,itemsep=2pt]
  \item \texttt{Controller::loadWidgets('main')} queries
        \texttt{NecoWidget::getRows()} which reads the
        \texttt{property} table.
  \item The widget routes are merged into \texttt{\$this->children} with
        string keys.
  \item \texttt{Controller::render()} iterates children, rendering each
        widget's HTML into \texttt{\$this->data[\$name . '\_code']}.
  \item \texttt{Controller::fetch()} requires the template, which
        includes \texttt{widgets-rows.tpl}.
  \item \texttt{widgets-rows.tpl} emits
        \texttt{\{\%widget\_name\%\}} tokens.
  \item \texttt{fetch()} substitutes the tokens with the widget HTML.
\end{enumerate}

This is the default mode. The admin has full control over which widgets
appear, in what order, in which columns, on which pages. No code changes
are needed.

\subsection{Manual Composition (Hardcoded Placeholders)}

In manual composition, a template author writes
\texttt{\{\%widget\_name\%\}} tokens directly in the \texttt{.tpl} file,
bypassing the admin widget manager. For this to work, the widget must
still be registered as a child of the controller (via
\texttt{loadWidgets()} or \texttt{addChild()}), so that its HTML is
rendered into \texttt{\$this->data[\$name . '\_code']}.

\begin{lstlisting}[caption={Manual widget placeholder hardcoding in a custom template.},label={lst:manual-hardcode}]
<!-- A custom landing page template that hardcodes widget positions -->
<div class="hero-section">
    {%hero_banner_widget%}
</div>

<div class="features-grid">
    <div class="col-4">{%product_overview_1%}</div>
    <div class="col-4">{%product_overview_2%}</div>
    <div class="col-4">{%product_overview_3%}</div>
</div>

<div class="content-section">
    {%richtext_widget_1%}
</div>

<div class="cta-section">
    {%contact_form_widget%}
</div>
\end{lstlisting}

For this to work, the controller must still call
\texttt{loadWidgets('main')} (or whatever position the widgets are
assigned to), so that the widget children are registered and their HTML
is rendered. The \texttt{widgets-rows.tpl} template is \emph{not}
included; instead, the custom template emits the placeholders directly
in the desired HTML structure.

\begin{findingbox}[Manual placeholders require widget children to be registered]
A \texttt{\{\%widget\_name\%\}} token in a template is only substituted
if the corresponding widget is registered as a child of the controller
(in \texttt{\$this->children} with the matching string key). This
happens automatically when \texttt{loadWidgets()} runs and the
\texttt{property} EAV table has a widget with that name for the current
position. If the admin has not configured a widget named
\texttt{hero\_banner\_widget} for the current page's position, the
placeholder token will remain in the HTML unsubstituted (it will appear
as literal text \texttt{\{\%hero\_banner\_widget\%\}} in the browser).

So manual composition requires coordination between the template author
(who writes the placeholder) and the admin (who must create a widget
instance with the matching name in the visual editor).
\end{findingbox}

\subsection{Hybrid Composition}

A template can combine both modes: use \texttt{widgets-rows.tpl} for
some positions (dynamic) and hardcode placeholders for others (manual).
For example, the \texttt{landing\_page\_necotienda.tpl} template
includes \texttt{widgets-common.tpl} (which uses
\texttt{widgets-rows.tpl} for dynamic composition) but also appends
hardcoded SVG sections after the widget-driven content:

\begin{lstlisting}[caption={Hybrid composition in \texttt{landing\_page\_necotienda.tpl}.},label={lst:hybrid}]
<?php echo $header; ?>
<?php $tpl = is_dir(...) ? $this->config->get('config_template') : "choroni"; ?>

<div id="contentContainer" class="tpl-page" nt-editable>
    <!-- Dynamic widget composition (via widgets-common.tpl) -->
    <?php include(DIR_TEMPLATE . $tpl . "/shared/widgets-common.tpl"); ?>
</div>

<!-- Hardcoded content (not widget-driven) -->
<div class="svg-container-01" style="width: 2918px; height: 1306px;">
    <?php include("landing_page_necotienda_svg_01.tpl"); ?>
</div>
<div class="svg-container-02" style="width: 2374px; height: 681px;">
    <?php include("landing_page_necotienda_svg_02.tpl"); ?>
</div>
\end{lstlisting}

This is the most flexible mode: the admin controls the main content area
via the visual editor, and the template author controls the surrounding
sections with hardcoded HTML.

% ============================================================
\section{The \texttt{only:} Position Prefix and Embedded Pages}\label{sec:only-prefix}
% ============================================================

\subsection{The \texttt{page\_embed.tpl} Layout}

When a page is embedded (via \texttt{ControllerContentPage::embed()}, as
called by the menu system's \texttt{submenu\_type = 'page\_id'}, v7),
the controller uses a simplified layout:

\begin{lstlisting}[caption={\texttt{page\_embed.tpl} --- the simplified embed layout.},label={lst:page-embed}]
<?php $tpl = ...; ?>
<div class="tpl-page-embed">
    <?php include(DIR_TEMPLATE . $tpl . "/shared/widgets-featured.tpl");?>

    <div class="row">
        <div class="large-12 medium-12 small-12">
            <div nt-editable>
                <?php $position = 'main'; ?>
                <?php include(DIR_TEMPLATE . $tpl . "/shared/widgets-rows.tpl");?>
            </div>
        </div>
    </div>

    <?php include(DIR_TEMPLATE . $tpl . "/shared/widgets-featured-footer.tpl");?>
</div>
\end{lstlisting}

No side columns, no breadcrumbs --- just the featured, main, and
featured-footer positions, all in a single full-width column.

\subsection{The \texttt{only:} Prefix}

When embedding a page, the controller uses the \texttt{only:} position
prefix:

\begin{lstlisting}[caption={The \texttt{only:} prefix in \texttt{ControllerContentPage::embed()}.},label={lst:only-prefix}]
$this->loadWidgets('only:featuredContent');
$this->loadWidgets('only:main');
$this->loadWidgets('only:featuredFooter');
\end{lstlisting}

The \texttt{only:} prefix tells \texttt{Controller::loadWidgets()} to
load \emph{only} the per-object widgets (those matching the session's
\texttt{object\_type}/\texttt{object\_id}), skipping the default widget
tree. This is because an embedded page should only show its own widgets,
not the entire page's widget tree.

In \texttt{Controller::loadWidgets()}, the \texttt{only:} prefix is
detected at line 510:

\begin{lstlisting}[caption={\texttt{Controller::loadWidgets()} --- the \texttt{only:} detection.},label={lst:only-detection}]
if (!strpos($position, 'only:')) {
    // Normal: load default tree + per-object tree
    $rows = $widgets->getRows($params, ...);
    // ... merge default rows ...
} else {
    // 'only:' prefix: skip the default tree
    // (the per-object query still runs below)
}
// Then, regardless of 'only:', the per-object query runs:
if ($session->has('object_type') || $session->has('object_id')) {
    $params['position'] = $position = str_replace('only:', '', $position);
    // ... query per-object widgets and merge ...
}
\end{lstlisting}

So \texttt{only:} skips the first query (the default widget tree) and
only runs the second query (the per-object widget tree). The result is
that an embedded page shows only its own widgets, not the page's default
widgets.

% ============================================================
\section{The CMS Lifecycle Trace}\label{sec:trace}
% ============================================================

\subsection{Visitor Views a Blog Post}

\begin{enumerate}[leftmargin=1.6em,itemsep=2pt]
  \item Request: \texttt{GET /articulos/my-first-post}
  \item SEO URL rewriter resolves \texttt{my-first-post} to
        \texttt{post\_id=42} via \texttt{url\_alias} table.
  \item Route: \texttt{content/post}.
  \item \texttt{ControllerContentPost::index()}:
        \begin{enumerate}[itemsep=2pt]
          \item Sets session: \texttt{object\_type = 'post'},
                \texttt{object\_id = 42},
                \texttt{landing\_page = 'content/post/index'}.
          \item Loads post: \texttt{modelPost->getById(42)} ---
                \texttt{SELECT * FROM post p LEFT JOIN description pd ON
                (pd.object\_id = p.post\_id AND pd.object\_type = 'post'
                AND pd.language\_id = <lang>)}.
          \item Checks customer group permission via
                \texttt{modelPost->getProperty(42, 'customer\_groups',
                'customer\_groups')}.
          \item Loads widgets: \texttt{loadWidgets('featuredContent')},
                \texttt{loadWidgets('main')},
                \texttt{loadWidgets('featuredFooter')}. Each runs two
                queries: default tree + per-object tree (filtered by
                \texttt{LIKE '\%object\_type=post\%'} and
                \texttt{LIKE '\%object\_id=42\%'}).
          \item Adds layout children: \texttt{common/header},
                \texttt{common/footer},
                \texttt{common/column\_left},
                \texttt{common/column\_right}.
          \item Template resolution: checks EAV
                \texttt{property(42, 'style', 'view')} for a per-post
                template. If empty, uses \texttt{config('default\_view\_post')}
                or falls back to \texttt{content/post.tpl}.
          \item \texttt{render(true)}:
                \begin{itemize}[itemsep=2pt]
                  \item Renders header, footer, column\_left,
                        column\_right (each a child controller).
                  \item Renders each widget child (each a module
                        controller).
                  \item Fetches \texttt{content/post.tpl}, which
                        includes \texttt{widgets-common.tpl}, which
                        includes \texttt{widgets-rows.tpl}, which emits
                        \texttt{\{\%widget\_name\%\}} tokens.
                  \item Substitutes tokens with widget HTML.
                  \item Post-processes HTML (strip comments, collapse
                        whitespace).
                  \item Returns the final HTML.
                \end{itemize}
          \item Cache write-back: \texttt{cache->set('html-post.42...',
                \$html)}.
          \item \texttt{response->setOutput(\$html)}.
        \end{enumerate}
\end{enumerate}

\subsection{Menu Embeds a Page}

\begin{enumerate}[leftmargin=1.6em,itemsep=2pt]
  \item The ``links'' widget renders a menu link with
        \texttt{submenu\_type = 'page\_id'},
        \texttt{page\_id = 7}.
  \item \texttt{ControllerModuleLinks::getLinks()} calls
        \texttt{content/page->embed(7)}.
  \item \texttt{ControllerContentPage::embed(7)}:
        \begin{enumerate}[itemsep=2pt]
          \item Sets session: \texttt{object\_type = 'page'},
                \texttt{object\_id = 7},
                \texttt{landing\_page = 'content/page/index'}.
          \item Loads widgets with \texttt{only:} prefix:
                \texttt{loadWidgets('only:featuredContent')},
                \texttt{loadWidgets('only:main')},
                \texttt{loadWidgets('only:featuredFooter')}. This loads
                \emph{only} the per-object widgets for page 7, skipping
                the default tree.
          \item Template: \texttt{content/page\_embed.tpl} (simplified
                layout, no side columns).
          \item \texttt{render(true)} --- renders the embedded page
                HTML.
          \item Returns the HTML to the links widget, which inserts it
                as the link's submenu content.
        \end{enumerate}
\end{enumerate}

% ============================================================
\section{Cross-Cutting Findings}\label{sec:findings}
% ============================================================

\subsection{Every CMS Page is a Widget Composition Surface}

The most important finding of v8 is that every CMS page --- whether a
blog post, a static page, or a category listing --- is a widget
composition surface. The same \texttt{loadWidgets()}/\texttt{addChild()}
/\texttt{render()} pipeline that drives the home page and product pages
also drives CMS pages. This means the admin can configure different
widgets for each individual post, page, or category --- a ``related
posts'' widget on one blog post, a ``contact form'' widget on a specific
page, a ``product list'' widget on a category --- all via the visual
editor, without code changes.

\subsection{The \texttt{widgets-common.tpl} Central Layout is a Refactoring}

The \texttt{widgets-common.tpl} fragment was introduced to centralise the
layout that was originally inlined in \texttt{home.tpl}. The home page
template still uses the old inline pattern (manually including each
fragment); newer templates (post, page, product, category) use the
centralised \texttt{widgets-common.tpl}. This is an evolutionary
refactoring scar --- the home page was never migrated.

\subsection{Manual Placeholders Require Admin Coordination}

A template author can hardcode \texttt{\{\%widget\_name\%\}} tokens in a
custom template, but this requires the admin to create a widget instance
with the matching name in the visual editor. If the admin doesn't, the
token appears as literal text in the browser. This coordination
requirement makes manual composition fragile --- a widget renamed or
deleted in the admin breaks the template.

\subsection{The \texttt{only:} Prefix is an Optimisation}

The \texttt{only:} prefix halves the number of
\texttt{NecoWidget::getRows()} queries for embedded pages (one query
instead of two). This is important because embedded pages are rendered
inside the links widget, which may render multiple embedded pages for a
single menu's submenu items. Without \texttt{only:}, each embedded page
would fire two queries; with \texttt{only:}, it fires one.

\subsection{Posts and Pages Share Everything But the Discriminator}

The post and page models are nearly identical: same table
(\texttt{post}), same primary key (\texttt{post\_id}), same fields, same
relations (\texttt{descriptions}, \texttt{stores}). The only difference
is \texttt{\$object\_type} (\texttt{'post'} vs \texttt{'page'}) and the
page model's \texttt{embed()} method. The \texttt{post\_type} column in
the database provides the discriminator. This is the same
single-table-with-discriminator pattern used by many CMS platforms
(WordPress's \texttt{post\_type}, Drupal's \texttt{bundle}).

% ============================================================
\section{Closing Observations}\label{sec:observations}
% ============================================================

\subsection{The CMS + Widget Integration is the Platform's Crown Jewel}

The integration of the CMS subsystem with the widget composition system
is the platform's crown jewel. It allows a non-technical admin to create
a blog post, assign widgets to it (a ``related posts'' carousel, a
``comments'' widget, a ``social share'' widget), choose a per-post
template (a standard blog layout or a custom landing page), and publish
--- all without touching a single line of code. The widget composition
system (v3) and the per-entity template override (v8) together form a
page-builder experience that rivals dedicated CMS platforms.

\subsection{The Manual Composition Mode is a Safety Valve}

The manual composition mode (hardcoding \texttt{\{\%widget\_name\%\}}
tokens) is a safety valve for cases where the visual editor's layout
capabilities are insufficient. A template author can create a custom
template with a fixed HTML structure and insert widget placeholders at
specific positions. This gives the admin control over the widget
\emph{content} (via the visual editor) while the template author
controls the widget \emph{layout} (via the hardcoded HTML).

\subsection{Final Note}

v8 confirms that the Necoyoad CMS subsystem is a widget-driven content
management system where every page is a composition surface. The shared
\texttt{post} table, the per-entity template override, the
\texttt{widgets-common.tpl} universal layout, and the
\texttt{\{\%widget\_name\%\}} placeholder substitution together form a
flexible page-builder architecture that supports both dynamic
(admin-configured) and manual (hardcoded) composition modes.

\appendix
\section{Glossary}\label{app:glossary}
\begin{table}[htbp]
\centering
\caption{Glossary of CMS and widget composition terms.}\label{tab:glossary-v8}
\footnotesize\sloppy
\begin{tabularx}{\columnwidth}{>{\raggedright\arraybackslash}p{3.5cm}X}
\toprule
\textbf{Term} & \textbf{Meaning} \\
\midrule
\texttt{post} table & Shared by blog posts and static pages, discriminated by \texttt{post\_type}. \\
\texttt{post\_type} & The discriminator: \texttt{'post'} for blog posts, \texttt{'page'} for static pages. \\
Dynamic composition & Widgets configured via the visual editor, stored in the \texttt{property} EAV table, rendered via \texttt{widgets-rows.tpl}. \\
Manual composition & \texttt{\{\%widget\_name\%\}} tokens hardcoded in a \texttt{.tpl} file. Requires matching widget instances in the admin. \\
Hybrid composition & A template that combines dynamic (via \texttt{widgets-common.tpl}) and hardcoded sections. \\
\texttt{widgets-common.tpl} & The universal layout fragment: featured + breadcrumbs + columns + featured-footer. \\
\texttt{widgets-rows.tpl} & The template that iterates the widget tree and emits \texttt{\{\%widget\_name\%\}} tokens. \\
\texttt{only:} prefix & A position prefix that loads only per-object widgets, skipping the default tree. Used by embedded pages. \\
Per-entity template override & EAV \texttt{property('style', 'view')} returns a template path; falls back to config default; falls back to \texttt{choroni/}. \\
\texttt{page\_embed.tpl} & Simplified layout for embedded pages (no side columns). Used by the menu system's \texttt{submenu\_type = 'page\_id'}. \\
\bottomrule
\end{tabularx}
\end{table}
\end{document}
